\begin{landscapepage}
	%----------------------------------------------------------------------
	\section*{Timeline}
	\label{sec:time}
	%----------------------------------------------------------------------

	If you already have a strong idea of the approach, and also some milestones to reach at a certain point, it makes sense to define a schedule so that you can check your progress against it later on. This can be done in a simple tabular format, or using a Gantt chart as shown below.

	\begin{center}
		\resizebox{.7\linewidth}{!}{%
			\begin{ganttchart}[
					hgrid,
					vgrid,
					bar height=.6,
					group/.append style={fill=black!30},
					bar/.append style={fill=black!10},
					milestone/.append style={fill=black},
					link/.style={->, thick},
					x unit=.55cm,
					y unit chart=.8cm,
					title height=.6,
					bar label font=\scriptsize,
					group label font=\scriptsize\bfseries,
					milestone label font=\scriptsize\itshape,
					bar label node/.append style={align=left, text width=4cm},
					group label node/.append style={align=left, text width=4cm}
				]{1}{14}

				% Year across all weeks
				\gantttitle{2026}{14} \\

				% Months with week shares (exemplary Jul/Aug/Sep)
				\gantttitle{Jul}{5}
				\gantttitle{Aug}{4}
				\gantttitle{Sep}{5} \\

				% Week numbers (1..12)
				\gantttitlelist{1,...,14}{1} \\

				% Group
				\ganttgroup[name=g1]{Phase A}{1}{6} \\

				% Tasks (bar)
				\ganttbar[name=task1]{Task 1}{1}{3} \\
				\ganttbar[name=task2]{Task 2}{3}{6} \\

				% Milestone
				\ganttmilestone[name=ms]{Milestone}{6} \\

				% Links
				\ganttlink{task1}{task2}
				\ganttlink{task2}{ms}

			\end{ganttchart}%
		}
	\end{center}
\end{landscapepage}
